\documentclass[11pt,a4paper]{article}

\usepackage{fontspec}
\usepackage[french]{babel}
\usepackage{geometry}
\geometry{margin=2.2cm}
\usepackage{graphicx}
\usepackage{booktabs}
\usepackage{array}
\usepackage{xcolor}
\usepackage{listings}
\usepackage{fancyhdr}
\usepackage{enumitem}
\usepackage{titlesec}
\usepackage{hyperref}
\usepackage{caption}
\usepackage{float}

\definecolor{bleu}{HTML}{1F4E79}
\definecolor{gris}{HTML}{F2F2F2}
\definecolor{vert}{HTML}{2E7D32}

\hypersetup{
  colorlinks=true,
  linkcolor=bleu,
  urlcolor=bleu,
  pdftitle={Rapport Examen TP DevOps - Biblio},
  pdfauthor={LACHHAB Mohammed}
}

\captionsetup{font=small,labelfont=bf}
\setlength{\parindent}{0pt}
\setlength{\parskip}{6pt}

\pagestyle{fancy}
\fancyhf{}
\lhead{\small Examen TP DevOps --- Biblio}
\rhead{\small LACHHAB Mohammed --- GINF2}
\cfoot{\thepage}
\renewcommand{\headrulewidth}{0.4pt}

\titleformat{\section}{\large\bfseries\color{bleu}}{\thesection}{0.6em}{}
\titleformat{\subsection}{\normalsize\bfseries\color{bleu}}{\thesubsection}{0.6em}{}

\lstset{
  basicstyle=\ttfamily\small,
  backgroundcolor=\color{gris},
  frame=single,
  rulecolor=\color{gray!40},
  breaklines=true,
  columns=fullflexible,
  keepspaces=true,
  showstringspaces=false,
  xleftmargin=4pt,
  xrightmargin=4pt,
  aboveskip=6pt,
  belowskip=6pt
}

\newcommand{\shot}[3]{%
  \begin{figure}[H]
    \centering
    \includegraphics[width=0.95\linewidth,height=0.46\textheight,keepaspectratio]{screenshots/#1}
    \caption{#2}
    \label{fig:#1}
  \end{figure}
}

\begin{document}

% ------------------------------------------------------------------
% Page de garde
% ------------------------------------------------------------------
\begin{titlepage}
  \centering
  \vspace*{1.2cm}

  {\large \textbf{Universite Abdelmalek Essaadi}}\\[0.3cm]
  {\large Ecole Nationale des Sciences Appliquees de Tanger}\\[0.2cm]
  {\large Departement Informatique}\\[1.4cm]

  {\Large \textbf{Examen TP DevOps}}\\[0.5cm]
  {\LARGE \textbf{Biblio}}\\[0.4cm]
  {\large GitFlow, GitHub Actions et conteneurisation Docker}\\[1.8cm]

  \begin{tabular}{rl}
    \textbf{Filiere :} & Genie Informatique (GINF2)\\
    \textbf{Realise par :} & LACHHAB Mohammed\\
    \textbf{Encadrants :} & Pr. Mohammed Ghailani\\
                         & Pr. Marouane Haddad\\
    \textbf{Depot :} & \href{https://github.com/edwinlabrinth/biblio-devops}{github.com/edwinlabrinth/biblio-devops}\\
  \end{tabular}

  \vfill
  {\large Annee universitaire 2026--2027}
\end{titlepage}

\tableofcontents
\newpage

% ------------------------------------------------------------------
\section{Contexte et objectifs}

Le projet \textbf{Biblio} est une application PHP en ligne de commande de gestion
de bibliotheque (livres, adherents, emprunts). Ce rapport presente le travail
individuel realise autour du code fourni : mise en place d'un depot
collaboratif GitFlow, automatisation de l'integration continue avec GitHub
Actions, conteneurisation Docker et operations de maintenance Git.

Le travail a ete mene seul. Les roles normalement repartis dans une equipe de
trois ont donc ete assures successivement par le meme etudiant. La publication
sur Docker Hub n'a pas ete activee : le workflow correspondant existe mais reste
conditionne, conformement au perimetre choisi.

\begin{table}[H]
  \centering
  \small
  \begin{tabular}{>{\bfseries}p{4.2cm}p{10.2cm}}
    \toprule
    Objectif & Resultat \\
    \midrule
    GitFlow & Branches \texttt{main}, \texttt{develop}, \texttt{feature/*}, \texttt{release/*}, \texttt{hotfix/*} \\
    Commits atomiques & Commits conventionnels, un sujet par commit \\
    Pull requests & 9 PR relues puis fusionnees avec un commit de merge \\
    Protection & \texttt{main} et \texttt{develop} protegees, force push interdit \\
    CI & Lint PHP, 39 tests, quality gate bloquant \\
    Docker & Image multi-stage, tests au build, utilisateur non-root \\
    Release & Tags SemVer \texttt{v1.0.0} et \texttt{v1.0.1} \\
    Docker Hub & Job prepare mais desactive (hors perimetre) \\
    \bottomrule
  \end{tabular}
  \caption{Synthese des objectifs et de leur etat.}
\end{table}

% ------------------------------------------------------------------
\section{Organisation GitFlow}

\subsection{Branches et commits}

Deux branches permanentes structurent le depot : \texttt{main} pour les
versions publiees et \texttt{develop} pour l'integration. Chaque fonctionnalite
est developpee dans une branche \texttt{feature/*}, puis fusionnee par une pull
request avec un commit de merge (pas de fast-forward). La figure
\ref{fig:01_git_graph.png} montre l'historique obtenu.

\shot{01_git_graph.png}{Historique GitFlow du depot (commande \texttt{git log --graph}).}{01_git_graph.png}

Les commits suivent la convention \texttt{type: description} :
\texttt{feat}, \texttt{fix}, \texttt{test}, \texttt{ci}, \texttt{build},
\texttt{docs}, \texttt{chore}. Le decoupage reste volontairement simple :
un commit par modele, un commit par harnais de tests, un commit par brique
d'infrastructure.

\subsection{Pull requests et protection}

Chaque integration passe par une PR avec description, checklist et relecture.
Les branches \texttt{main} et \texttt{develop} sont protegees :

\begin{itemize}[leftmargin=1.2cm]
  \item pull request obligatoire avant fusion ;
  \item statut \texttt{Tests \& Quality Gate} obligatoire et a jour ;
  \item force push et suppression de branche interdits ;
  \item fusion par commit de merge uniquement (squash et rebase desactives).
\end{itemize}

\shot{06_pull_requests.png}{Pull requests fermees et fusionnees sur GitHub.}{06_pull_requests.png}

Le travail etant individuel, la relecture croisee a ete remplacee par une
auto-relecture explicite : un commentaire de revue est depose sur chaque PR
avant la fusion, et la CI doit etre verte. Le depot est visible a l'adresse
\url{https://github.com/edwinlabrinth/biblio-devops}.

% ------------------------------------------------------------------
\section{Integration et livraison continues}

Le workflow \texttt{CI} se declenche sur les push et les pull requests vers
\texttt{main} et \texttt{develop}. Il contient deux jobs principaux :

\begin{table}[H]
  \centering
  \small
  \begin{tabular}{>{\bfseries}p{4.2cm}p{9.8cm}}
    \toprule
    Job & Etapes \\
    \midrule
    \texttt{Tests \& Quality Gate} & installation de PHP 8.4, verification de syntaxe,
      execution des 39 assertions, quality gate bloquante, publication du rapport \\
    \texttt{Build \& Smoke Test Docker} & construction de l'image (tests au build),
      verification de l'UID non-root, execution ephemere en lecture seule \\
    \bottomrule
  \end{tabular}
  \caption{Jobs du workflow d'integration continue.}
\end{table}

Le quality gate est bloquant : si le script de tests retourne un code non nul,
la PR ne peut pas etre fusionnee. La figure \ref{fig:07_actions_ci.png} montre
les executions vertes sur la branche \texttt{develop}.

\shot{07_actions_ci.png}{Executions du workflow CI sur GitHub Actions.}{07_actions_ci.png}

Le workflow \texttt{Release} est declenche par un tag \texttt{vX.Y.Z}. Il
rejoue les tests, construit l'image et calcule automatiquement les tags SemVer
\texttt{X.Y.Z}, \texttt{X.Y}, \texttt{X} et \texttt{latest}. La figure
\ref{fig:08_release_run.png} montre l'execution reussie de la version 1.0.1.

\shot{08_release_run.png}{Workflow de release declenche par le tag \texttt{v1.0.1}.}{08_release_run.png}

\shot{09_tags.png}{Tags SemVer publies sur le depot.}{09_tags.png}

% ------------------------------------------------------------------
\section{Conteneurisation Docker}

L'image est construite en trois etapes : \texttt{base}, \texttt{test} et
\texttt{runtime}. L'etape \texttt{test} execute \texttt{php tests/run.php}
pendant le build ; l'image finale herite de cette etape, donc elle ne peut pas
etre produite si un test echoue. Le conteneur s'execute avec un utilisateur
non privilegie et peut tourner en lecture seule.

\begin{lstlisting}
FROM php:8.4-cli-alpine AS base
WORKDIR /app
COPY autoload.php index.php VERSION ./
COPY src ./src
COPY tests ./tests
RUN addgroup -S biblio && adduser -S -G biblio biblio
USER biblio

FROM base AS test
RUN php tests/run.php

FROM test AS runtime
ENTRYPOINT ["php", "index.php"]
\end{lstlisting}

\shot{03_docker_build.png}{Construction de l'image : les tests sont executes au build.}{03_docker_build.png}

L'execution est ephemere et isolee :

\begin{lstlisting}
docker run --rm -i --read-only --tmpfs /tmp \
  --security-opt no-new-privileges:true biblio:1.0.1
\end{lstlisting}

\shot{04_docker_run.png}{Execution isolee du conteneur avec une entree utilisateur.}{04_docker_run.png}

Le fichier \texttt{docker-compose.yml} reprend les memes options
(\texttt{read\_only}, \texttt{tmpfs}, \texttt{no-new-privileges}) pour un usage
local. Les tests locaux restent executables sans Docker avec
\texttt{php tests/run.php} (figure \ref{fig:02_tests.png}).

\shot{02_tests.png}{Suite de tests : 39 assertions, aucune erreur.}{02_tests.png}

% ------------------------------------------------------------------
\section{Maintenance Git}

Les operations de maintenance ont ete documentees dans
\texttt{docs/maintenance.md}. L'inspection de l'historique et le reflog ont ete
utilises pour retrouver les commits et verifier l'etat des branches
(figure \ref{fig:12_maintenance.png}).

\shot{12_maintenance.png}{Inspection de l'historique, attribution des lignes et reflog.}{12_maintenance.png}

\begin{table}[H]
  \centering
  \small
  \begin{tabular}{>{\bfseries}p{4.8cm}p{9.2cm}}
    \toprule
    Besoin & Commande / principe \\
    \midrule
    Inspection & \texttt{git log --graph}, \texttt{git show}, \texttt{git blame} \\
    Recherche de regression & \texttt{git bisect start / good / bad / reset} \\
    Remise a zero controlee & \texttt{git reset --soft}, \texttt{--mixed}, recuperation par \texttt{reflog} \\
    Compensation & \texttt{git revert <commit>}, sans reecrire l'historique \\
    \bottomrule
  \end{tabular}
  \caption{Commandes de maintenance prevues et documentees.}
\end{table}

La recherche dichotomique n'a pas ete executee sur une regression reelle du
domaine : la suite de tests et la CI ont detecte les erreurs au moment des PR.
La procedure reste documentee et applicable telle quelle. Le hotfix 1.0.1
illustre la compensation d'un defaut decouvert apres la release 1.0.0 : sortie
infinie du menu lorsque l'entree standard est fermee.

% ------------------------------------------------------------------
\section{Bilan}

\begin{itemize}[leftmargin=1.2cm]
  \item Depot GitFlow complet, branches protegees et historique lisible.
  \item 9 pull requests relues, CI verte avant chaque fusion.
  \item 39 assertions, quality gate bloquant et rapport de tests archive.
  \item Image Docker multi-stage, non-root, testee au build et executee en
        mode ephemere.
  \item Deux versions SemVer publiees : \texttt{v1.0.0} et \texttt{v1.0.1}.
  \item Publication Docker Hub volontairement non activee.
\end{itemize}

Le principal point d'adaptation est le travail individuel : la revue croisee
et la repartition des roles ont ete remplacees par une auto-relecture et une
execution sequentielle des taches. Le resultat reste tracable dans les PR et
dans l'historique Git.

\section*{Annexe : commandes utiles}
\addcontentsline{toc}{section}{Annexe : commandes utiles}

\begin{lstlisting}
# tests
php tests/run.php

# image et execution
docker build -t biblio:1.0.1 .
docker run --rm -it --read-only --tmpfs /tmp biblio:1.0.1

# etat du depot
git log --graph --oneline --decorate --all
git tag -n
git status -sb
\end{lstlisting}

\end{document}
